\documentclass[11pt]{article}
\usepackage[a4paper,margin=1in]{geometry}
\usepackage{fontspec}
\usepackage[boldfont=false]{xeCJK}
\setCJKmainfont{FandolSong-Regular.otf}
\usepackage{amsmath}
\usepackage{amssymb}
\usepackage{array}
\usepackage{booktabs}
\usepackage{graphicx}
\usepackage{adjustbox}
\usepackage{enumitem}
\setlist[itemize]{topsep=2pt,partopsep=0pt,parsep=0pt,itemsep=2pt,leftmargin=1.4em}
\setlist[enumerate]{topsep=2pt,partopsep=0pt,parsep=0pt,itemsep=2pt,leftmargin=1.6em}
\usepackage{parskip}
\usepackage{fvextra}
\fvset{breaklines=true,breakanywhere=true,fontsize=\small}
\setlength{\parindent}{0pt}

\begin{document}

\begin{center}
  {\LARGE\bfseries Excretion in humans}\\[0.35em]
  {\large Igcse Biology}
\end{center}
\vspace{0.6em}

\section*{Excretion}

\textbf{Excretion} 排泄 is the removal of the waste products of metabolism, and of substances the body has in excess. Two main wastes are removed:

\begin{itemize}
\item \textbf{carbon dioxide} 二氧化碳 — made during respiration, and breathed out through the lungs.

\item \textbf{urea} 尿素 — made in the \textbf{liver} 肝脏, and removed by the kidneys together with excess water and \textbf{ions} 离子.

\end{itemize}

\section*{The urinary system}

The \textbf{kidneys} 肾脏 clean the blood and make \textbf{urine} 尿液. The urine then passes through:

\begin{itemize}
\item the \textbf{ureters} 输尿管 — tubes from the kidneys to the bladder.

\item the \textbf{bladder} 膀胱 — stores the urine.

\item the \textbf{urethra} 尿道 — carries the urine out of the body.

\end{itemize}

\section*{Inside the kidney (Supplement)}

A kidney has two regions: an outer \textbf{cortex} 皮质 and an inner \textbf{medulla} 髓质.

\subsection*{The nephron}

Each kidney holds millions of tiny tubes called \textbf{nephrons} 肾单位. A nephron cleans the blood in two steps:

\begin{enumerate}
\item \textbf{Filtration} 过滤: at the \textbf{glomerulus} 肾小球 (a knot of capillaries), high pressure forces water, \textbf{glucose} 葡萄糖, urea and ions out of the blood and into the nephron.

\item \textbf{Reabsorption} 重吸收: as the liquid flows along the nephron, the blood takes back \textbf{all} of the glucose, \textbf{some} of the ions, and \textbf{most} of the water.

\end{enumerate}

What is left behind — urea, excess water and excess ions — becomes the urine.

\section*{The liver, amino acids and urea (Supplement)}

The liver carries out the \textbf{assimilation} 同化 of \textbf{amino acids} 氨基酸: it joins them together into the \textbf{proteins} 蛋白质 that the body needs.

The body cannot store extra amino acids. The liver breaks the excess down by \textbf{deamination} 脱氨基作用 — it removes the \textbf{nitrogen} 氮-containing part of each amino acid. This part is then turned into urea.

Urea must be excreted because it is \textbf{toxic} 有毒 (poisonous) if it is allowed to build up in the blood.

\section*{Exam tips}

\begin{itemize}
\item Excretion removes metabolic wastes. Carbon dioxide leaves through the lungs; urea, excess water and excess ions leave through the kidneys.

\item Urine path: kidney → ureter → bladder → urethra.

\item Nephron: the glomerulus \textbf{filters} the blood; the nephron \textbf{reabsorbs} all the glucose, some ions and most of the water. What is left is urine.

\item Urea is made in the \textbf{liver} from excess amino acids by \textbf{deamination}.

\item We must excrete urea because it is \textbf{toxic} to the body.

\end{itemize}


\end{document}
